% =====================================================================
%  BAB 1 — PENDAHULUAN
% =====================================================================

\chapter{Pendahuluan}

\section{Latar Belakang}
\label{sec:1.1}

[Uraikan latar belakang masalah secara komprehensif. Mulai dari kondisi umum, identifikasi
kesenjangan (\emph{gap}) yang ada, hingga urgensi penelitian ini dilakukan. Dukung dengan
data dan referensi empiris. Panjang: 3--5 paragraf.]

[Paragraf pertama: konteks besar / kondisi terkini di bidang ini.]

[Paragraf kedua: masalah spesifik yang belum terpecahkan / keterbatasan pendekatan yang ada.]

[Paragraf ketiga: solusi yang diusulkan secara ringkas dan mengapa pendekatan ini relevan.]

\section{Rumusan Masalah}
\label{sec:1.2}

Berdasarkan latar belakang di atas, penelitian ini merumuskan permasalahan sebagai berikut:

\begin{enumerate}
  \item Bagaimana [rumusan masalah pertama]?
  \item Bagaimana [rumusan masalah kedua]?
  \item Seberapa [rumusan masalah ketiga, biasanya terkait kinerja/evaluasi]?
\end{enumerate}

\section{Tujuan Penelitian}
\label{sec:1.3}

Penelitian ini bertujuan untuk:

\begin{enumerate}
  \item [Tujuan pertama, sesuai dengan rumusan masalah 1.]
  \item [Tujuan kedua, sesuai dengan rumusan masalah 2.]
  \item [Tujuan ketiga, sesuai dengan rumusan masalah 3.]
\end{enumerate}

\section{Batasan Masalah}
\label{sec:1.4}

Agar penelitian terfokus, ditetapkan batasan-batasan sebagai berikut:

\begin{enumerate}
  \item [Batasan 1: lingkup teknologi, dataset, atau platform yang digunakan.]
  \item [Batasan 2: jenis eksperimen atau skenario yang dicakup.]
  \item [Batasan 3: asumsi yang disederhanakan dari kondisi nyata.]
\end{enumerate}

\section{Manfaat Penelitian}
\label{sec:1.5}

\textbf{Manfaat Teoritis:}
[Jelaskan kontribusi terhadap pengembangan ilmu pengetahuan, misalnya: memperkaya literatur
di bidang X, menyediakan kerangka referensi untuk penelitian sejenis, dsb.]

\medskip
\textbf{Manfaat Praktis:}
[Jelaskan manfaat langsung bagi praktisi atau industri, misalnya: membantu tim keamanan
dalam Y, menyederhanakan proses Z, dsb.]

\section{Sistematika Penulisan}
\label{sec:1.6}

Tesis ini disusun dalam tujuh bab dengan sistematika sebagai berikut.

\textbf{Bab~1 Pendahuluan} menyajikan latar belakang, rumusan masalah, tujuan, batasan,
dan manfaat penelitian.

\textbf{Bab~2 Tinjauan Pustaka} membahas teori-teori dasar dan penelitian terdahulu yang
relevan sebagai landasan penelitian ini.

\textbf{Bab~3 Metodologi Penelitian} menguraikan desain penelitian, metode pengumpulan data,
dan teknik analisis yang digunakan.

\textbf{Bab~4 Perancangan Sistem} mendeskripsikan arsitektur dan desain komponen sistem yang
dibangun.

\textbf{Bab~5 Implementasi} menjelaskan realisasi sistem berdasarkan perancangan pada Bab~4.

\textbf{Bab~6 Pengujian dan Analisis} menyajikan hasil eksperimen beserta analisis dan
pembahasan secara mendalam.

\textbf{Bab~7 Kesimpulan dan Saran} merangkum temuan utama penelitian dan mengusulkan arah
penelitian lanjutan.
